% English follow-up insertion; root integrates into the main manuscript.
% Raw records and limitations: audit/persistence_boundary_review.md.
% No new bibliography entries are introduced by this insertion.
The persistence follow-up examined the application boundary and the execution
environment without changing the original Store or repeating the 32 native
replays.  Source inspection found synchronous SQLite DELETE/EXTRA transactions,
commit-before-receipt behavior, and no delayed application writer in the
reviewed path.  The environment did expose an active workspace synchronization
service and a writable rsync route to the workspace.  Both the workspace and
\texttt{/tmp} were on the same overlay mount with
\texttt{fsync=volatile}; hence the previously observed location contrast was
not a contrast between filesystem mounts.  As specified by the Linux overlay
documentation, this mount mode omits synchronization to its upper filesystem.
Successful SQLite or explicit synchronization calls therefore do not establish
physical or power-loss durability here.  Neither this mount property nor the
identified synchronization route proves the cause of the historical prefix
regressions.

Two bounded experiments compared plain fixed-length files with explicit
synchronization, ordinary SQLite DELETE/EXTRA databases, and the unchanged
Store using a pure counter adapter, at both locations.  The paced experiment
covered 24 files with 14 updates and a 12-second final observation interval;
the burst experiment covered 36 files with 14 updates within a single
wall-clock second for each file and a 10-second final interval.  All 60 live
controls preserved their final bytes, logical state and inode identity,
including subsequent checks in separate executions.  Measurements retained
file descriptors to distinguish the original inode from the current pathname,
and recorded hashes, device/inode identifiers, nanosecond timestamps, explicit
close/reopen results and filesystem notifications.  A separate synthetic
positive control deliberately replaced each of three final pathnames with its
own earlier prefix: all three replacements were detected while the retained
final inodes remained intact.  This validates the detector, but supplies no
historical causal attribution.  The original residues and replay failures
remain preserved; neither live experiment reproduced them.

The operational workflow consequently keeps live databases outside automatic
file mirroring and publishes only closed, consistently captured snapshots.
In this environment \texttt{/tmp} separates live work from the observed
workspace synchronization route, but remains transient and shares the same
volatile mount; persistent deployment requires a suitable storage durability
contract.  A published snapshot is sealed by its byte hash and by a final
journal head/event anchor retained separately, verified after reopening and
again after transfer.  The sealed artifact is then treated as immutable.
This external anchor is essential because an internally consistent journal
may still be a valid earlier prefix.  Hashes and anchors provide integrity
checks, not authentication unless independently protected.
